\section{Implementation}

\subsection{Repository Organization}

The prototype is implemented in Python. Runtime code lives under
\texttt{src/}, command-line entry points under \texttt{src/cli/}, evaluation
drivers under \texttt{benchmarks/}, scenario definitions under
\texttt{benchmarks/scenarios/}, and DIAL Chat under \texttt{src/ui/}. The
repository is organized around the main system responsibilities:

\begin{itemize}
  \item \texttt{node.py} and \texttt{models.py} define the node runtime and
  shared data structures;
  \item the service modules implement discovery, profile exchange,
  recommendation requests, and query forwarding;
  \item \texttt{routing\_service.py} and \texttt{routing\_policy.py} implement
  capability matching, scoring, and node selection;
  \item \texttt{node\_api\_service.py} and \texttt{web\_app.py} expose the node
  API and DIAL Chat;
  \item \texttt{benchmarks/run\_benchmark.py} and
  \texttt{benchmarks/scenarios/} drive the local evaluation workloads.
\end{itemize}

Other files support this core path with environment loading, logging, backend
adapters, health checks, protocol helpers, and small development utilities.
The command-line tools and scripts are used to start nodes, launch local
networks, inspect discovery behavior, and run diagnostic queries. The main
user-facing interface is DIAL Chat, which lives under \texttt{src/ui/} and is
described later in this chapter.

The public source code is available at \url{https://github.com/SpicyPaper/DIAL}.

\subsection{Node Runtime and Startup}

A single node can be started with \texttt{src/cli/run\_node.py}, while a local
network of several nodes can be launched with
\texttt{scripts/start\_network.sh}. Both read their configuration from the
project \texttt{.env} file; full setup instructions are kept in the project
\texttt{README}.

When a node starts, it creates a libp2p host, derives its peer identity,
initializes its profile, connects to bootstrap peers, starts the DHT service,
and advertises its top capabilities. The network script launches several nodes
with different ports, model names, and profiles. Runtime state is written under
\texttt{.runtime/}.

The provided local setup advertises IPv6 loopback multiaddresses
(\texttt{/ip6/::1/...}) because this was the local address mode used by the
test machine when several nodes were run on one host. This is a deployment
choice, not a routing assumption. libp2p
multiaddresses can describe different network addresses and transport paths
\citep{libp2paddressing}, so supporting IPv4, public IPv6, or relay-based
addresses such as circuit relay paths \citep{libp2pcircuitrelay} would mainly
require extending the address-selection mode, launch configuration, and
connectivity tests.

The concrete \texttt{NodeProfile} dataclass contains the fields needed by this
runtime: \texttt{peer\_id}, reachable \texttt{addresses}, \texttt{model\_name},
\texttt{advertised\_capabilities}, full \texttt{capability\_scores},
\texttt{is\_available}, and \texttt{timestamp\_ms}. Architecturally, these are
the profile data described in Chapter 4; here they become concrete fields used
by discovery, caching, and routing.

Capabilities are set when the node starts. A node receives configured
capabilities or role hints, and the local capability scorer turns them into a
deterministic profile: usually one strong primary capability and weaker
secondary capabilities. These scores are routing metadata, not measurements of
the model's real skill, because the project focuses on discovery and routing
rather than benchmarking the LLM itself.

\subsection{LLM Backends and Prompts}

The \texttt{LocalAgent} abstraction supports several model execution modes:

\begin{itemize}
  \item \textbf{dummy}: deterministic placeholder responses for tests and
  controlled routing runs;
  \item \textbf{Ollama}: local inference through an Ollama server, recommended
  for demos;
  \item \textbf{AIaaS}: EPFL OpenAI-compatible inference API;
  \item \textbf{local Transformers}: a Hugging Face Transformers backend for
  fully local model loading.
\end{itemize}

This modularity matters because a decentralized AI network should not assume
one runtime. A node may call a model from a local GPU, CPU runtime,
institutional API, or specialized inference engine.

The local network is intentionally easy to run. Normal usage defaults to Ollama
with the configured Qwen3 4B model \citep{ollama}; controlled runs can also use
a dummy backend when the goal is to exercise routing without generation
variance.

In real runs, answer quality depends on the chosen backend model and on the
system prompt used for the node. This is expected: prompting choices can change
LLM behavior and robustness \citep{schulhoff2025promptreport,zhu2024promptrobust}.
The prompts used in this prototype are basic role prompts, not optimized prompts
or a contribution of the project.

For practical reasons, the provided local network script normally launches
several nodes against the same configured LLM backend. This is a
resource-friendly demo setup, not an architectural restriction: each node process
has its own peer identity, advertised profile, routing state, and system prompt,
and custom launches can use different backend configurations. Backend calls are
stateless unless conversation context is explicitly sent, so one node does not
automatically remember another node's previous conversation.

\subsection{LLM-Based Capability Classification}

In normal interactive use, the entry node uses a classifier LLM to map the user
prompt to capability scores. The supported capability set includes general,
math, programming, writing, summarization, research, planning, and creative
tasks. The classifier is guided by a dedicated system prompt and returns a small
JSON object of scores, which makes classification an explicit routing step
rather than an unstructured chat response.

The classifier output is not a final node choice. It is a compact description of
what the query needs, and the router keeps at most the top three capabilities
for discovery and scoring. Its quality depends on the selected classifier model
and prompt, so controlled evaluation runs can bypass this LLM step by injecting
\texttt{required\_capabilities} directly. This lets Chapter 7 evaluate the
network and routing behavior without also measuring classifier quality.

\subsection{Routing and Query Execution}

The routing implementation connects classification, discovery, and execution.
For each query, \texttt{routing\_service.py} normalizes the required
capabilities, searches the DHT for peers advertising the relevant capabilities,
fetches full profiles through the profile protocol, stores them in the peer
registry, checks candidate liveness, and computes the routing score. The first
candidate pool contains the local node and available matching registry profiles
after this DHT refresh. If no candidate in that pool passes the threshold, the
recommendation protocol asks eligible remote candidates for additional peer IDs.
Those recommended peers are fetched and scored with the same routing logic.

Once a candidate is selected, \texttt{query\_service.py} performs the execution
step. If the selected candidate is the local node, the node calls its local
agent directly. If the selected candidate is remote, the query is forwarded
through the query protocol and the selected peer calls its own bound LLM
backend. The response is then returned to the entry node and back to the
client.

\subsection{HTTP API and DIAL Chat}

Each node exposes a small HTTP API through FastAPI \citep{fastapi}:

\begin{itemize}
  \item \texttt{POST /api/query}: submit a query to the node;
  \item \texttt{GET /api/query/progress/<query\_id>}: read routing progress
  events;
  \item \texttt{GET /api/profile}: inspect the node profile.
\end{itemize}

DIAL Chat reads available entry-node APIs from runtime configuration and
presents a conversational interface. A user can choose an entry node, send a
prompt, keep a conversation, and see how much recent context will be sent with
the next request. The context meter makes clear that long conversations remain
in the interface, while only the recent context window is included in the
network request.

\begin{figure}[H]
\centering
\includegraphics[width=\textwidth]{DIALChatUserQuery.png}
\caption{DIAL Chat main interface. The user selects an entry node, sends a query, and sees the recent conversation context that will be included in the next network request.}
\label{fig:dial-chat-main}
\end{figure}

DIAL Chat also exposes routing progress events: classification, DHT lookup,
candidate evaluation, recommendation, selection, forwarding, and answer
generation. This makes the prototype easier to inspect because the user can see
which node answered and why the routing path took time.

\begin{figure}[H]
\centering
\includegraphics[width=\textwidth]{DIALChatWithNetworkDetails.png}
\caption{DIAL Chat network details panel. The interface shows timing information, required capabilities, candidate evaluation, and the selected peer that generated the answer.}
\label{fig:dial-chat-network-details}
\end{figure}

\subsection{Evaluation Harness}

The evaluation runner starts local networks, injects scenario parameters, sends
queries, records routing traces, and writes summaries. It can vary network size,
latency, dropout, workload mix, and routing policy. Chapter 7 describes the
actual scenarios and results.

The harness is designed to evaluate the network and routing layer, not the
underlying LLMs. During benchmark runs, both the answer backend and classifier
backend are set to \texttt{dummy}. Benchmark queries also provide
\texttt{required\_capabilities} directly in the API request, so the classifier is
not called: the scenario defines the capability scores that the router should
use. The selected node then returns a dummy answer instead of an LLM-generated
one. This keeps the measured results focused on discovery, scoring,
recommendation, forwarding, and failure handling. In normal interactive use, by
contrast, the required capabilities come from the classifier LLM and the selected
node calls its configured LLM backend.

\subsection{Implementation Limitations}

The current prototype is the result of a semester project, not a production
network. Some mechanisms are therefore simplified deliberately so the project
can focus on building and measuring the decentralized discovery and routing
layer:

\begin{itemize}
  \item capability scores are configured or simulated routing metadata, not
  measurements of real model skill;
  \item Gossipsub is implemented only as an optional profile-update notification
  path; normal local runs rely on DHT lookup and direct peer protocols;
  \item nodes have libp2p peer identities, but profiles are not signed and the
  prototype does not implement a trust mechanism for verifying advertised
  capabilities;
  \item unreachable peers are handled after failed connection or response
  attempts: the peer is marked unavailable, excluded, and routing is retried,
  but this availability information is not yet shared with the rest of the
  network;
  \item evaluation runs on a local machine with IPv6 loopback addresses and
  does not validate IPv4/public IPv6 addressing, NAT traversal, relay behavior,
  wide-area network behavior, or real provider heterogeneity;
  \item malicious behavior, false capability advertisements, and generated
  answer correctness are outside the evaluated scope;
  \item incentive, payment, and reputation systems are future work.
\end{itemize}

These limitations clarify the result rather than weaken it. As Chapter 7 shows,
DIAL is a working prototype and gives encouraging routing results in the tested
local scenarios. Several missing pieces, including stronger verification,
reputation, incentives, privacy, and richer failure handling, are discussed as
future improvements in Chapter 8.
